
\documentclass[conference]{IEEEtran}

\usepackage{amsmath,amssymb}
\usepackage{booktabs}
\usepackage{graphicx}
\usepackage{url}
\usepackage{array}
\usepackage{xcolor}

\begin{document}

\title{NeuroRank 2.0: An Empirical Study of
Retrieval-Aware Learning-to-Rank and
Query-Dependent Model Selection for
E-Commerce Search}

\author{
\IEEEauthorblockN{Bhukya Kasi}
\IEEEauthorblockA{
Indian Institute of Technology Dharwad\\
Karnataka, India\\
is24bm028@iitdh.ac.in
}
}

\maketitle

\begin{abstract}
E-commerce search systems commonly combine
lexical retrieval, neural semantic retrieval,
and supervised learning-to-rank. However,
ranking models trained on predefined labeled
candidates may generalize poorly when deployed
on candidates generated by a retrieval pipeline.
This paper presents NeuroRank 2.0, a modular
experimental search platform integrating BM25,
MiniLM sentence embeddings, FAISS retrieval,
reciprocal rank fusion, LambdaMART reranking,
and query-dependent adaptive model selection.

The study investigates the behavior of these
methods under a retrieval-aware experimental
protocol using the Amazon Shopping Queries
ESCI dataset. A query-disjoint research
partition includes 400 training queries,
100 validation queries, and 100 held-out
test queries. The fixed experimental catalog
contains 13,083 products.

On the held-out test queries, reciprocal rank
fusion achieved the highest observed
NDCG@10 of 0.5445, followed by semantic
retrieval at 0.5407, adaptive ranking at
0.5325, LambdaMART at 0.5088, and BM25
at 0.4981. The paired difference between
fusion and semantic retrieval was small
(0.00386), with a 95\% bootstrap confidence
interval spanning zero. These results
indicate that more complex ranking strategies
do not necessarily improve end-to-end
relevance in sparse-judgment, retrieval-
dependent settings. The platform also
provides FastAPI inference and an
interactive Streamlit dashboard.
\end{abstract}

\begin{IEEEkeywords}
Information Retrieval, E-Commerce Search,
Learning-to-Rank, LambdaMART, BM25,
Sentence Transformers, FAISS,
Reciprocal Rank Fusion, Model Selection
\end{IEEEkeywords}

\section{Introduction}

Product search requires matching short,
often ambiguous customer queries with
relevant items from a large catalog.
Traditional lexical methods are effective
when query and document vocabularies
overlap, but may struggle with semantic
paraphrases. Dense neural retrieval
addresses some vocabulary mismatches
by mapping queries and products into
a shared vector space.

Modern retrieval systems frequently
apply an additional supervised ranking
model to the retrieved candidates.
However, a learning-to-rank model may
encounter candidate distributions during
deployment that differ from those seen
during supervised training. Such
differences can affect feature values,
label availability, and ranking quality.

NeuroRank 2.0 was developed to examine
these practical challenges through an
integrated retrieval and ranking system.
The platform combines lexical and neural
retrievers with supervised and adaptive
ranking approaches.

The main contributions of this study are:

\begin{itemize}
\item An implemented multi-stage product
search platform combining BM25, MiniLM,
FAISS, reciprocal rank fusion, and
LambdaMART.

\item A retrieval-aware supervised ranking
protocol that uses explicitly judged
query-product pairs selected from
retrieval-generated candidate pools.

\item An experimental query-dependent
selector that chooses between semantic
ranking and LambdaMART.

\item A held-out comparison of five
ranking strategies, supported by
feature ablation, robustness diagnostics,
and paired statistical evaluation.
\end{itemize}

The work is primarily an empirical and
engineering contribution. It does not
claim the invention of BM25, LambdaMART,
or reciprocal rank fusion.

\section{Related Work}

\subsection{Lexical and Neural Retrieval}

BM25 is a widely used probabilistic
term-weighting approach to document
retrieval \cite{robertson}.
Its scoring function rewards matching
query terms while accounting for term
frequency and document length.

Sentence-BERT introduced an efficient
approach to obtaining semantically
meaningful sentence representations
\cite{reimers}. Compact MiniLM models
can be used to produce dense query and
product representations.

FAISS provides efficient similarity
search over dense vectors
\cite{johnson}.

\subsection{Learning-to-Rank}

Supervised learning-to-rank uses relevance
judgments to learn a ranking function
over query-document features. LambdaMART
combines boosted decision trees with
ranking-oriented optimization
\cite{burges}. LightGBM provides an
efficient implementation of gradient-
boosted trees \cite{ke}.

\subsection{Rank Fusion}

Reciprocal rank fusion (RRF) combines
ranked lists without requiring their
raw similarity scores to be calibrated
to a common scale \cite{cormack}.
This makes it a useful baseline when
different retrieval models produce
incomparable scores.

\subsection{Product Search Benchmarks}

The Amazon Shopping Queries dataset
contains query-product relevance
judgments under the ESCI scheme:
Exact, Substitute, Complement, and
Irrelevant \cite{reddy}. It enables
empirical comparisons of product
ranking approaches.

\section{System Architecture}

NeuroRank 2.0 follows a two-stage
retrieval-and-ranking architecture.

Given a query $q$ and product catalog
$\mathcal{D}$, lexical and semantic
retrievers independently generate
ranked candidate lists.

The candidate union is

\begin{equation}
\mathcal{C}(q)=
\mathcal{C}_{B}(q)
\cup
\mathcal{C}_{S}(q),
\end{equation}

where $\mathcal{C}_{B}$ denotes BM25
candidates and $\mathcal{C}_{S}$
denotes semantic candidates.

Each retriever contributes at most
100 candidates. The union contains
at most 200 unique products.

\subsection{BM25 Retrieval}

For a query $q$ and document $d$,
BM25 computes

\begin{equation}
\begin{split}
s_B(q,d)=
\sum_{t\in q}
\mathrm{IDF}(t)
\frac{
f(t,d)(k_1+1)
}{
f(t,d)+
k_1\left(
1-b+
b\frac{|d|}{\mathrm{avgdl}}
\right)
}.
\end{split}
\end{equation}

Here $f(t,d)$ is the frequency
of term $t$ in document $d$,
$\mathrm{IDF}(t)$ is its inverse
document frequency, and $k_1$
and $b$ control normalization.

\subsection{Dense Semantic Retrieval}

A MiniLM encoder $f_{\theta}$
maps queries and products into
dense vectors.

The semantic score is

\begin{equation}
s_S(q,d)=
\frac{
f_{\theta}(q)^\top
f_{\theta}(d)
}{
\|f_{\theta}(q)\|
\|f_{\theta}(d)\|
}.
\end{equation}

Normalized vectors are indexed using
FAISS inner-product search.

\subsection{Reciprocal Rank Fusion}

For candidate $d$, the fusion score is

\begin{equation}
s_{\mathrm{RRF}}(d)=
\sum_{r\in\{B,S\}}
\frac{
\mathbf{1}[d\in\mathcal{C}_{r}]
}{
60+\mathrm{rank}_{r}(d)
}.
\end{equation}

The indicator is zero when a
candidate is absent from a
retriever's list.

\subsection{LambdaMART Reranking}

For each retrieved query-product pair,
the supervised ranker uses the
eight-dimensional feature vector

\begin{equation}
\mathbf{x}(q,d)=
[
s_B,s_S,o,e,r_B,r_S,l_d,l_q
]^\top,
\end{equation}

where $o$ is query-title token
overlap, $e$ indicates exact
query phrase occurrence in the
title, $r_B$ and $r_S$ are reciprocal
rank features, and $l_d,l_q$
are title and query token lengths.

A LambdaMART model predicts a
ranking score

\begin{equation}
s_L(q,d)=F_{\mathrm{LGBM}}
(\mathbf{x}(q,d)).
\end{equation}

ESCI labels are mapped to
graded relevance values

\begin{equation}
I\mapsto0,\quad
C\mapsto1,\quad
S\mapsto2,\quad
E\mapsto3.
\end{equation}

Corresponding ranking gains are
$[0,1,3,7]$.

\subsection{Query-Dependent Selection}

An experimental logistic regression
selector receives query-only statistics,
including word count, character
count, lexical diversity, average
word length, and the number of tokens
containing digits.

It predicts

\begin{equation}
p(q)=P(L>S\mid\mathbf{z}(q)),
\end{equation}

where $L>S$ denotes that LambdaMART
has higher judged-candidate
NDCG@10 than semantic ranking
on a validation query.

The deployed experimental rule is

\begin{equation}
s_A(q,d)=
\begin{cases}
s_L(q,d), & p(q)>0.5,\\
s_S(q,d), & \text{otherwise}.
\end{cases}
\end{equation}

This is a query-dependent selector,
not a verified uncertainty-calibrated
reliability estimator.

\section{Experimental Methodology}

\subsection{Dataset and Partitioning}

The experiments use a US-language
subset of the Amazon ESCI dataset.

A previous pilot study used a
2,936-product catalog and 30 test
queries. Those queries were examined
during development and were not used
as the final independent evaluation.

A separate research partition
contains 400 training queries,
100 validation queries, and
100 held-out test queries.
Previously used query identifiers
and normalized query text were
excluded during construction.

The resulting fixed catalog
contains 13,083 products.

The catalog includes text from
products associated with the three
splits. Consequently, the study
tests generalization to unseen
queries within a fixed catalog,
rather than generalization to
entirely unseen products.

\subsection{Retrieval-Aware Training}

For a training query, BM25 and
MiniLM each retrieve at most
100 products.

Only candidate pairs with
available ESCI relevance labels
are used for supervised training.

Queries lacking adequate
judged label variation are omitted
from LambdaMART fitting.

Unjudged retrieved products are
not automatically labeled as
irrelevant during supervised training.

\subsection{Evaluation Measures}

Ranking quality is evaluated with
normalized discounted cumulative gain:

\begin{equation}
\mathrm{DCG}@k=
\sum_{i=1}^{k}
\frac{
2^{rel_i}-1
}{
\log_2(i+1)
}.
\end{equation}

\begin{equation}
\mathrm{NDCG}@k=
\frac{
\mathrm{DCG}@k
}{
\mathrm{IDCG}@k
}.
\end{equation}

Recall at rank $k$ is

\begin{equation}
\mathrm{Recall}@k=
\frac{
|\mathcal{R}_q\cap
\mathrm{Top}_k(q)|
}{
|\mathcal{R}_q|
},
\end{equation}

where $\mathcal{R}_q$ is the set
of known relevant products.

For end-to-end NDCG, unjudged
retrieved products are assigned
zero gain. This is an evaluation
assumption, not a claim that such
products are truly irrelevant.

\subsection{Statistical Protocol}

The held-out evaluation uses
100 query-level observations.

Paired 95\% percentile bootstrap
confidence intervals use
10,000 resamples.

Two-sided paired sign-flip
permutation tests use
20,000 resamples.

Holm correction is applied to
five preselected NDCG@10
comparisons.

\section{Results and Discussion}

\subsection{Held-Out End-to-End Evaluation}

Table~\ref{tab:main} presents the
held-out results.

\begin{table}[t]
\centering
\caption{Held-Out Ranking Performance
(100 Queries; 13,083 Products)}
\label{tab:main}
\footnotesize
\setlength{\tabcolsep}{3.0pt}
\begin{tabular}{lccc}
\toprule
Method &
NDCG@10 &
Recall@10 &
Recall@100 \\
\midrule
BM25 & 0.4981 & 0.3589 & 0.6432\\
MiniLM & 0.5407 & 0.3965 & 0.7538\\
RRF & \textbf{0.5445} &
\textbf{0.4018} &
\textbf{0.7682}\\
LambdaMART & 0.5088 & 0.3692 & 0.7536\\
Adaptive & 0.5325 & 0.3885 & 0.7549\\
\bottomrule
\end{tabular}
\end{table}

RRF achieved the largest observed
NDCG@10 and recall values among
the evaluated approaches.

MiniLM was a competitive baseline,
with only a 0.0038 NDCG@10 gap
relative to RRF.

LambdaMART and adaptive ranking
did not surpass the stronger
retrieval baselines.

The candidate-pool recall was
approximately 0.7946. Thus,
some known relevant products
were absent before reranking.

\subsection{Validation Versus Test Behavior}

On 97 usable validation queries,
judged-candidate NDCG@10 was:

\begin{itemize}
\item MiniLM: 0.7795;
\item LambdaMART: 0.7708;
\item adaptive selector: 0.7786.
\end{itemize}

These values are not directly
comparable with end-to-end test
NDCG because the judged-candidate
evaluation excludes unjudged
retrieved products.

The adaptive selector chose
LambdaMART for 10 out of
100 held-out queries.

\subsection{Statistical Comparisons}

The mean paired difference
in NDCG@10 between RRF and
MiniLM was approximately
$0.00386$.

Its 95\% bootstrap confidence
interval was approximately

\begin{equation}
[-0.0327,\;0.0392].
\end{equation}

The Holm-adjusted permutation
$p$-value was 1.0.

Therefore, the experiment does
not establish a statistically
significant RRF improvement
over MiniLM.

For RRF versus BM25, the
observed paired mean difference
was approximately $0.0464$,
with a 95\% bootstrap interval
of approximately

\begin{equation}
[0.0223,\;0.0724].
\end{equation}

The latter result provides
stronger evidence of an advantage
for rank fusion over lexical
retrieval in this sample.

\subsection{Feature Ablation and Diagnostics}

Feature-group ablations were
performed by removing lexical,
semantic, text-matching, or
length-related features from
LambdaMART and retraining with
the same overall configuration.

The ablation outputs are recorded
in the project's Phase 11D
experiment artifacts. Numerical
ablation comparisons should be
inserted into the final manuscript
from the corresponding saved file
rather than reconstructed from
memory.

Earlier diagnostic experiments
also revealed differences between
features of predefined labeled
candidates and features of live
retrieved products.

These observations motivate
further study of training-serving
candidate distribution mismatch.
They do not independently prove
that distribution mismatch caused
the test performance gap.

\section{Engineering Implementation}

The experimental system uses
Python, Pandas, NumPy,
sentence-transformers, FAISS,
rank-bm25, LightGBM,
and scikit-learn.

FastAPI exposes a search endpoint
with selectable ranking modes.
A Streamlit interface enables
interactive product search and
side-by-side comparisons.

A separate local benchmark over
the earlier 2,936-product pilot
catalog measured mean latency
of 34.82 ms and p95 latency
of 40.00 ms for 100 warmed
sequential requests.

These latency measurements are
not estimates for the larger
13,083-product catalog and do
not constitute a concurrent
production load test.

\section{Limitations}

Several limitations should be
considered when interpreting
the findings.

First, relevance judgments are
sparse. Treating unjudged
products as zero-gain items
can underestimate actual
ranking relevance.

Second, only 100 final test
queries are evaluated.
The catalog is a sampled
subset of the full product
collection.

Third, the selector is based
on a small set of query-text
statistics. Its output should
not be interpreted as
calibrated reliability.

Fourth, the retrieval-aware
supervised model learns only
from retrieved products for
which human labels are available.

Finally, the experiment is
an observational comparison of
ranking pipelines. It does not
causally isolate every possible
source of training-serving shift.

\section{Conclusion}

This paper presented NeuroRank 2.0,
an integrated experimental platform
for lexical search, neural retrieval,
rank fusion, supervised reranking,
and query-dependent model selection.

On a 100-query held-out benchmark,
RRF achieved the highest observed
NDCG@10 of 0.5445, but its small
advantage over MiniLM was not
statistically established.

The supervised and adaptive methods
did not outperform the strongest
retrieval baselines in the final
end-to-end evaluation.

These results highlight the
importance of assessing ranking
methods under their actual
candidate-generation process,
rather than relying only on
predefined judged-candidate
evaluation.

Future work includes broader
candidate annotation, more diverse
query sets, better calibration,
and retrieval-aware feature learning.

\begin{thebibliography}{00}

\bibitem{robertson}
S. Robertson and H. Zaragoza,
``The probabilistic relevance
framework: BM25 and beyond,''
\emph{Foundations and Trends in
Information Retrieval},
vol. 3, no. 4,
pp. 333--389, 2009.

\bibitem{reimers}
N. Reimers and I. Gurevych,
``Sentence-BERT: Sentence
embeddings using Siamese
BERT-networks,''
in \emph{Proc. EMNLP-IJCNLP},
2019, pp. 3982--3992.

\bibitem{johnson}
J. Johnson, M. Douze, and H. Jegou,
``Billion-scale similarity
search with GPUs,''
\emph{IEEE Trans. Big Data},
vol. 7, no. 3,
pp. 535--547, 2021.

\bibitem{burges}
C. J. C. Burges,
``From RankNet to LambdaRank
to LambdaMART: An overview,''
Microsoft Research,
Tech. Rep. MSR-TR-2010-82, 2010.

\bibitem{ke}
G. Ke et al.,
``LightGBM: A highly efficient
gradient boosting decision tree,''
in \emph{Advances in Neural
Information Processing Systems},
2017.

\bibitem{cormack}
G. V. Cormack, C. L. A. Clarke,
and S. Buettcher,
``Reciprocal rank fusion
outperforms Condorcet and
individual rank learning methods,''
in \emph{Proc. ACM SIGIR},
2009, pp. 758--759.

\bibitem{reddy}
C. K. Reddy et al.,
``Shopping Queries Dataset:
A Large-Scale ESCI Benchmark
for Improving Product Search,''
arXiv:2206.06588, 2022.

\end{thebibliography}

\end{document}
